% 03_methodology -- written by Claude Code (D-045); revised after the peer review (Prompt 16: M2, M6, m1, m3, m4, m8,
% m11, m12, m16). Settings are quoted from configs/protocol_v1.yaml (frozen); every count comes from generated/numbers.tex.
\section{Methodology}\label{sec:method}

\subsection{Study design}
We ran a paired, within-item comparison: every test item was answered by the same model under two decoding
strategies, and each item contributes one correct-or-wrong outcome per strategy. The study answers three questions,
fixed before any model output on test items existed: \textbf{RQ1} (confirmatory): does the majority answer of 10
sampled runs have higher accuracy than the greedy answer? \textbf{RQ2} (descriptive): how consistent are the
answers across the 10 sampled runs? \textbf{RQ3}: does agreement between runs flag answers that are more likely to
be correct? Figure~\ref{fig:flow} summarises the pipeline.

\begin{figure}[t]
\centering
\begin{tikzpicture}[
  node distance=3.5mm and 4mm,
  box/.style={draw, rounded corners=2pt, align=center, font=\scriptsize, inner sep=2.5pt, text width=#1},
  box/.default=27mm,
  arr/.style={->, >=stealth, semithick}]
\node[box] (medqa) {MedQA, USMLE 4-option\\\medqaN{} items};
\node[box, right=of medqa] (pool) {Oncology keyword pool\\\poolN{} items};
\node[box, right=of pool] (screen) {AI screening M / O / X;\\blind AI re-screen of \screenN{}};
\node[box, right=of screen] (elig) {Eligible M\,+\,O: \eligibleN;\\image / bare-letter: $-$\excludedN};
\node[box, below=of elig] (split) {\finalN{} items, seeded split:\\debug \debugN{} $\cdot$ test \nTest};
\node[box, left=of split] (debug) {Debug set: prompt, parser and token limit fixed; protocol frozen};
\node[box, left=of debug] (model) {Qwen2.5-7B-Instruct, FP16,\\vLLM on 2\,$\times$\,T4};
\node[box, left=of model] (arms) {Greedy: $T$ = 0, 1 run\\Sampled: $T$ = 0.7, 10 runs};
\node[box, below=of arms] (parse) {Parse ``Answer: X'';\\majority vote of 10};
\node[box, right=of parse, text width=58mm] (eval) {RQ1 accuracy difference (exact McNemar, paired bootstrap)
  $\cdot$ RQ2 consistency (unanimity, Fleiss' $\kappa$) $\cdot$ RQ3 agreement as a flag (AUROC)};
\node[box, right=of eval, text width=27mm, dashed] (expl) {Exploratory (v2): greedy rerun, $T$ = 0.3 / 1.0, contamination probe};
\draw[arr] (medqa) -- (pool);
\draw[arr] (pool) -- (screen);
\draw[arr] (screen) -- (elig);
\draw[arr] (elig) -- (split);
\draw[arr] (split) -- (debug);
\draw[arr] (debug) -- (model);
\draw[arr] (model) -- (arms);
\draw[arr] (arms) -- (parse);
\draw[arr] (parse) -- (eval);
\draw[arr, dashed] (eval) -- (expl);
\end{tikzpicture}
\Description{A flowchart in three rows. Top row: MedQA items, oncology keyword pool, AI screening with a blind
re-screen, eligible items after exclusions. Middle row, right to left: the seeded split into debug and test sets,
the debug set used to fix and freeze the protocol, the model, and the two decoding arms. Bottom row: answer parsing
and majority vote, evaluation of the three research questions, and the exploratory runs.}
\caption{Study pipeline. Item selection (top row, left to right), protocol freeze and generation (middle row,
right to left), evaluation (bottom row). No model output on test items existed before the protocol was frozen.}
\label{fig:flow}
\end{figure}

\subsection{Items}
Items come from the English, four-option USMLE part of MedQA~\cite{jin2021medqa} on Hugging Face~\cite{medqa_hf}
(\medqaN{} items, pinned commit). A keyword rule (term list in the repository's \texttt{count\_oncology\_items.py}) kept
\poolN{} items with an oncology term in the vignette or the correct answer. Since the pool was drawn from all of MedQA,
the test set mixes items from its train (\nOriginTrain), dev (\nOriginDev) and test (\nOriginTest) splits, which
matters for contamination (Section~\ref{sec:limitations}). An AI screener (Claude Code) coded each item against written
rules: \emph{M}, managing a known or suspected neoplasm or a direct complication; \emph{O}, a neoplasm tested otherwise
(diagnosis, pathology, mechanism, risk, screening, adverse effects of cancer therapy); \emph{X}, incidental. With no
clinical checker available, a second, blind instance of the same model re-coded a random 20\% sample (\screenN{} items)
from the full text; they agreed on \screenAgree{} (Cohen's $\kappa$ = \screenKappa), and the first screener adjudicated
the \adjudicatedN{} disagreements, changing \adjudicatedChanged{} codes (\screenChangedPct\%, our estimate of the
miscoding rate among the \unsampledN{} unsampled items). Of \eligibleN{} eligible items (\codeM{} M, \codeO{} O),
\excludedN{} referring to an image or with bare-letter options were removed; a seeded split left \debugN{} items for
developing the prompt and parser and \nTest{} test items (\nM{} M, \nO{} O). An \emph{objective subset} of \nObjective{}
test items, whose correct answer itself contains an oncology term, depends less on screening judgement.

\subsection{Model, prompt and decoding}
We used Qwen2.5-7B-Instruct~\cite{qwen2024qwen25,qwen25_7b_instruct} at a pinned revision, in 16-bit floating
point, with its own chat template. The system message was \emph{``You are a careful medical expert. Answer the
multiple-choice question.''}; the user message gave the question, the options A--D in the dataset's order, and the
instruction to think briefly, step by step, and end with a line \emph{``Answer: X''}. \emph{Greedy decoding} used
temperature 0 and one output per item; \emph{sampling} drew 10 outputs per item at temperature 0.7, top-$p$ 0.8 and
top-$k$ 20, the model's own defaults. Both used a repetition penalty of 1.05 and at most 512 output tokens; a truncated
output is scored as wrong, which costs the single greedy answer more than the vote, so we also report a lenient
analysis. Seeds were a fixed base plus the item's index; the first 50 test items were run greedily twice.

\subsection{Answer extraction and majority vote}
The answer is the last ``Answer: X'' in the output (Markdown bold allowed), or else the last ``answer is X''; only
a single capital A--D not followed by another letter counts. An output without one is \emph{unparseable} and is
scored as wrong. The \emph{majority vote} is the most frequent letter among the parsed runs; unparseable runs do not
vote, a tie goes to the tied letter that appeared first, and an item with no parsed run has no vote. The \emph{modal
share} is the number of runs giving the voted letter, divided by 10.

\subsection{Outcomes and statistical analysis}
\paragraph{RQ1.} The primary outcome is the accuracy of the majority vote minus that of greedy decoding on the same
items, tested with the exact two-sided McNemar test~\cite{mcnemar1947} on the discordant items ($b$: vote right,
greedy wrong; $c$: the reverse), with a 95\% paired percentile bootstrap CI (10{,}000 resamples). The decision rule
was fixed in advance: \emph{supported} if $p < 0.05$ \emph{and} the lower bound is above 0; \emph{practically
meaningful} if the difference is at least 3 points. A power calculation for the then-planned 684 items gave a minimum
detectable difference of 3.7--4.3 points at a nominal 80\% power (77--78\% in simulation with the exact test) for
12--16\% discordant items. We also report the mean accuracy of a single sampled run and of a majority over the first
$k$ runs ($k$ = 1, 3, 5, 7, 9, 10).
\paragraph{RQ2 and RQ3.} Consistency is described (without a test) by the share of unanimous items, the mean modal
share, Fleiss' $\kappa$~\cite{fleiss1971} over the categories A--D and unparseable, and the pairwise flip rate (the
share of an item's 45 run pairs that disagree), with bootstrap CIs (2{,}000 resamples for $\kappa$ and the AUROC,
10{,}000 otherwise). For RQ3, the AUROC of the modal share for a correct vote is reported; H3 is \emph{supported} if its
lower bound exceeds 0.5, a deliberately modest bar. We also compare accuracy at $\geq$9/10 and $<$9/10 agreement and
plot accuracy against coverage.
\paragraph{Sensitivity.} RQ1 is repeated with unparseable greedy items dropped, on the objective subset, and in the M
and O subgroups (M is descriptive only). Discordant items and items wrong in all 11 generations are analysed. Only
RQ1 is confirmatory; secondary analyses are not corrected for multiple comparisons.

\subsection{Protocol freeze, amendment and post hoc analyses}
The protocol and its decision rules were frozen and committed to the project's repository (commit \texttt{7890add})
before any model output on test items existed, and acceptance criteria for the greedy baseline were written before
the test run. The AI assistant that designed the protocol had read the test items and keys during screening, but no
model output. After validating the baseline, and before looking at the voting comparison, we added an exploratory
amendment as a new protocol version: a greedy rerun of all test items, 10 sampled runs at $T$ = 0.3 and 1.0, and a
contamination probe in which the first wrong option of 100 random test items was hidden and the model asked to
reproduce its exact text (Wilson interval~\cite{wilson1927}). Analyses added after the results were seen, most of them
at the request of a reviewer, are labelled \emph{post hoc}.
